% vip-ribs-mvi-redux.tex — map of iOS architectures beyond MVC/MVVM/VIPER: Clean Swift (VIP),
% RIBs (Uber), MVI, Redux/ReSwift, TCA and the SwiftUI "MV" approach. One block per pattern —
% its loop drawn at column width, its mechanism beside it — and a who-owns-state table.
% Sources (checked 2026-10-05):
%   clean-swift.com/clean-swift-ios-architecture · github.com/Clean-Swift/CleanStore (ListOrders scene;
%     re-checked 2026-10-11: 6 files — ViewController, Interactor, Presenter, Router, Worker, Models)
%   github.com/uber/RIBs (README, wiki) · github.com/uber/RIBs-iOS RIBs/Classes/Interactor.swift
%   uber.com/blog/new-rider-app-architecture (2016-12-20)
%   staltz.com/unidirectional-user-interface-architectures.html (2015-08-22)
%   hannesdorfmann.com/android/mosby3-mvi-1 · redux.js.org (three principles; history of Redux)
%   github.com/ReSwift/ReSwift (README, releases) · github.com/pointfreeco/swift-composable-architecture
%   developer.apple.com/documentation/swiftui/managing-model-data-in-your-app
%   objc.io/issues/13-architecture/viper (Gilbert, Stoll, 2014-06) · redux.js.org/faq/organizing-state
%   developer.android.com/topic/architecture/ui-layer/events
% @labels: area=design kind=architecture platform=ios discussion=5
% @tags: vip, clean-swift, ribs, mvi, redux, reswift, mv-pattern, viper, interactor, unidirectional-data-flow, tca, observation
\documentclass[8pt]{extarticle}
\usepackage{printup-sheet}
\usepackage{array}

\lstdefinelanguage{SwiftSheet}{
  morekeywords={protocol,class,final,struct,enum,func,var,let,weak,init,
    if,else,return,guard,self,nil,try,await,async,private,some,is,
    true,false,case,switch,default,static,in,AnyObject},
  sensitive=true, morecomment=[l]{//}, morestring=[b]",
  literate={->}{{\hbox{-}\hbox{>}}}2 {==}{{\hbox{=}\hbox{=}}}2 {??}{{\hbox{?}\hbox{?}}}2}

% Every loop is drawn for ONE 80 mm column, with its boxes at the body's size (8 pt) and its
% labels at 7 pt: four loops in one row across the page needed 5–6 pt labels (owner,
% 2026-10-11: "rearange diagrams the way we do not need to make text smaller").
\tikzset{
  rb/.style={box, font=\normalsize, minimum width=17mm, minimum height=7.5mm, inner sep=2.5pt},
  vw/.style={rb, draw=sheetBlue, fill=sheetBlue!8},
  lg/.style={rb, draw=sheetGreen!60!black, fill=sheetGreen!12},
  pr/.style={rb, draw=sheetOrange, fill=sheetOrange!10},
  rt/.style={rb, draw=sheetBrown, fill=sheetBrown!10},
  gy/.style={rb, draw=sheetGrey, fill=black!4},
  sub/.style={font=\small},
  lbl/.style={font=\small, text=black!75, inner sep=1.5pt, align=center},
  note/.style={font=\small, text=sheetGreen!50!black, align=center},
  tn/.style={draw=sheetGreen!60!black, fill=sheetGreen!10, rounded corners=2pt, font=\normalsize,
             inner sep=2pt, align=center, minimum width=15mm, minimum height=7mm},
}

\begin{document}
\pdfonly{\folioslug{design/vip-ribs-mvi-redux}}

\sheettitle{VIP · RIBs · MVI · Redux · MV — the architecture map}{design · folio}

\oneliner{Every iOS architecture answers three questions: \textbf{who owns the state},
\textbf{which way does data flow}, and \textbf{where do side effects live}. \textbf{VIP} (Clean
Swift) makes each screen a one-way \emph{cycle} of three objects; \textbf{RIBs} (Uber) drive the app
by a \emph{tree of business logic}, not of views; \textbf{MVI} and \textbf{Redux} keep one
immutable state changed only by a reducer; \textbf{TCA} is Redux made composable; \textbf{MV} says
SwiftUI views plus \texttt{@Observable} models are already enough.}

% No "Data flow" column: each pattern's loop below draws it (say a fact once).
\vspace{2pt}
{\small\setlength{\tabcolsep}{3pt}\renewcommand{\arraystretch}{1.05}
\noindent\begin{tabular}{@{}>{\raggedright\arraybackslash}p{20mm}>{\raggedright\arraybackslash}p{36mm}>{\raggedright\arraybackslash}p{36mm}>{\raggedright\arraybackslash}p{28mm}>{\raggedright\arraybackslash}p{37mm}@{}}
\toprule
\textbf{Pattern} & \textbf{State lives in} & \textbf{Testability} & \textbf{Boilerplate} & \textbf{Fits} \\
\midrule
VIP & Interactor (+ \texttt{DataStore}) & high: spy the next role & high: 6 files per scene & UIKit apps, per-scene rigour \\
RIBs & each RIB's Interactor & high: I/R/B behind protocols & very high, codegen & huge apps, hundreds of engineers \\
MVI & one immutable screen State & high: pure reducer & medium & complex screens, Android parity \\
Redux / ReSwift & one global Store & high reducers; effects harder & medium–high & shared app-wide state \\
TCA & Store tree of features & very high: \texttt{TestStore} & high & big SwiftUI teams \\
MV & \texttt{@Observable} models & models yes, view logic no & low & SwiftUI small–mid apps \\
\bottomrule
\end{tabular}}

\begin{multicols}{2}
\raggedright

% ── VIP ────────────────────────────────────────────────────────────────────────────────────
\section{VIP (Clean Swift)}
\noindent\begin{tikzpicture}[sheet]
  \path[use as bounding box] (0,-0.62) rectangle (8,2.08);
  \node[vw] (vc) at (4.0,1.7) {ViewController};
  \node[pr] (pe) at (1.15,0.65) {Presenter};
  \node[lg] (in) at (6.85,0.65) {Interactor};
  \draw[flow] (vc.east) -| node[lbl, right, pos=0.75]{Request} (in.north);
  \draw[flow] (in) -- node[lbl, below]{Response} (pe);
  \draw[hot] (pe.north) |- node[lbl, left, pos=0.3]{ViewModel} node[lbl, above, pos=0.78, text=sheetOrange]{weak} (vc.west);
  \node[rt] (ro) at (1.15,-0.3) {Router \textnormal{\small owned by VC}};
  \node[gy] (wo) at (6.85,-0.3) {Worker \textnormal{\small API, DB}};
  \draw[flow] (in) -- (wo);
  \node[note] at (4.0,-0.3) {one-way cycle,\\one per scene};
\end{tikzpicture}

Each role is a protocol: VC = \texttt{DisplayLogic}, Interactor = \texttt{BusinessLogic} (holds the
\texttt{DataStore}), Presenter = \texttt{PresentationLogic} (formats the \emph{ViewModel}, calls
\texttt{display…}). \textbf{Router} (\texttt{RoutingLogic} + \texttt{DataPassing}) navigates,
passing data via the \texttt{DataStore}. Models nest per use case
(\texttt{ListOrders.FetchOrders.Request}). Older templates: \texttt{…Input}/\texttt{…Output}
protocols, wired by a \textbf{Configurator}.

\section{One VIP use case}
\begin{lstlisting}[language=SwiftSheet]
enum ListOrders { enum Fetch {          // per use case
  struct Request {}
  struct Response { let orders: [Order] }
  struct ViewModel { let rows: [String] } } }
protocol ListOrdersDisplayLogic: AnyObject {
  func display(_ vm: ListOrders.Fetch.ViewModel) }
final class ListOrdersInteractor {      // BusinessLogic
  var presenter: ListOrdersPresenter?; let worker = OrdersWorker()
  func fetch(_ r: ListOrders.Fetch.Request) {
    worker.load { self.presenter?.present(.init(orders: $0)) } } }
final class ListOrdersPresenter {       // PresentationLogic
  weak var viewController: ListOrdersDisplayLogic?  // no cycle
  func present(_ r: ListOrders.Fetch.Response) {
    viewController?.display(.init(rows: r.orders.map(\.title))) } }
\end{lstlisting}

% ── MVI ────────────────────────────────────────────────────────────────────────────────────
\section{MVI}
\noindent\begin{tikzpicture}[sheet]
  \path[use as bounding box] (0,-0.5) rectangle (8,1.85);
  \node[vw] (mv) at (4.0,1.5) {View \textnormal{\small\texttt{render(state)}}};
  \node[lg] (mm) at (1.3,0.4) {Model\\[-1pt]\textnormal{\small actions $\to$ state}};
  \node[pr] (mi) at (6.7,0.4) {Intent\\[-1pt]\textnormal{\small events $\to$ actions}};
  \draw[flow] (mv.east) -| node[lbl, right, pos=0.75]{events} (mi.north);
  \draw[flow] (mi) -- node[lbl, above]{actions} (mm);
  \draw[flow] (mm.north) |- node[lbl, left, pos=0.3]{state} (mv.west);
  \node[note] at (4.0,-0.3) {one immutable state per screen · RxJS Observables in Cycle.js};
\end{tikzpicture}
% (no text under the loop: it carries all of MVI — the View is a pure render(state))

\section{MV and TCA}
\textbf{MV}: no per-screen ViewModel; logic lives in \texttt{@Observable} models (iOS 17, macOS 14).
A view updates only when a property its \texttt{body} reads changes. Own with \texttt{@State},
share with \texttt{.environment(model)}, read with \texttt{@Environment(Model.self)}, bind with
\texttt{@Bindable}. The \texttt{Observable} protocol alone observes nothing; the
\texttt{@Observable} macro does. \textbf{TCA}: State · Action · Reducer · Store; effects and dependencies are
injectable, and \texttt{TestStore} asserts every state change and effect.

% ── RIBs ───────────────────────────────────────────────────────────────────────────────────
\section{RIBs (Uber)}
\noindent\begin{tikzpicture}[sheet]
  \path[use as bounding box] (0,-0.9) rectangle (8,2.2);
  \node[tn] (root) at (4.0,1.85) {Root · R · I · V};
  \node[tn] (lo) at (1.4,0.88) {LoggedOut\\R · I · V};
  \node[tn, draw=sheetOrange, fill=sheetOrange!10, dashed] (li) at (5.6,0.88) {LoggedIn\\R · I \textbf{(no view)}};
  \node[tn] (og) at (4.45,-0.12) {OffGame\\R · I · V};
  \node[tn] (tt) at (6.95,-0.12) {Game\\R · I · V};
  \draw[flow] (root) -- (lo); \draw[flow] (root) -- (li);
  \draw[flow] (li) -- (og); \draw[flow] (li) -- (tt);
  \draw[hot] (tt.north east) to[bend right=50] (li.east);
  \node[lbl, align=center] at (1.4,-0.12) {Router attaches /\\detaches children};
  \node[note] at (4.0,-0.72) {tree of \emph{logic}; views follow · \textcolor{sheetOrange}{orange = Listener (up)}};
\end{tikzpicture}

\textbf{R}outer · \textbf{I}nteractor · \textbf{B}uilder (+ Component for dependencies, optional
Presenter and View). The Interactor holds business logic and Rx subscriptions, alive between
\texttt{didBecomeActive()} and \texttt{willResignActive()}; bind with
\texttt{disposeOnDeactivate(interactor:)}. Builder (factory + DI) creates the unit and is the only
part that knows the DI system. Up: a \texttt{Listener} the
parent implements; down: Rx streams or \texttt{build()} arguments.

\textbf{Why:} Uber's rider app was MVC grown to a few hundred engineers, one
\texttt{RequestViewController} from 300 to 3,000+ lines. RIBs make \textbf{business logic, not the
view tree, drive routing}.

% ── Redux / ReSwift ────────────────────────────────────────────────────────────────────────
\section{Redux / ReSwift}
\noindent\begin{tikzpicture}[sheet]
  \path[use as bounding box] (0,-0.82) rectangle (8,2.12);
  \node[vw] (rv) at (4.0,1.75) {View \textnormal{\small\texttt{StoreSubscriber}}};
  \node[gy] (mw) at (6.8,0.98) {Middleware};
  \node[lg] (rr) at (6.8,-0.12) {Reducer \textnormal{\small pure}};
  \node[rt] (rs) at (1.4,0.35) {Store\\[-1pt]\textnormal{\small\textbf{one} app State}};
  % the label rides ABOVE the horizontal run: left of the drop it sat on the View box's corner
  \draw[flow] (rv.east) -| node[lbl, above, pos=0.3]{\texttt{dispatch(a)}} (mw.north);
  \draw[flow] (mw) -- (rr);
  \draw[flow] (rr.west) -- (rs.south east);
  \draw[flow] (rs.north) |- node[lbl, left, pos=0.3, align=right]{\texttt{newState}\\\texttt{(state:)}} (rv.west);
  \node[note] at (4.0,-0.65) {single source of truth · async: thunks (ReSwift-Thunk)};
\end{tikzpicture}

Redux's 3 principles: one store · state is read-only, changed only by an action · pure reducers
(state + action $\to$ new state). ReSwift:
\texttt{Store<AppState>(reducer:state:)}; \texttt{subscribe(self) \{ \$0.select(…) \}} narrows
updates, skipping repeats of an \texttt{Equatable} substate.

\section{Traps}
\begin{itemize}
  \trap{``VIP is VIPER renamed.'' VIPER (2014): Presenter \emph{between} View and Interactor,
        both ways; Entities stay in the Interactor. VIP: a one-way cycle.}
  \trap{RIBs' Router is not ``just navigation'': it attaches child RIBs from Interactor
        output; screens follow.}
  \trap{Redux: subscribers hear every change unless they \texttt{select} a slice; local UI
        state (\texttt{isDropdownOpen}) stays out of the store.}
  \trap{A one-off message kept in state shows again on every re-render until the view
        reports it shown and the state clears it.}
\end{itemize}

\end{multicols}

\noindent{\footnotesize\color{sheetGrey}\textit{Related:} MVC / MVP / MVVM · Coordinator, Repository,
DI, Clean / VIPER · The Composable Architecture (TCA) · unidirectional state machines · Modularization \& SPM}

\end{document}
